\chapter[\emj{💥}~Hashing a Fondo: Colisiones e Implementación]{\emj{💥}~Hashing a Fondo: Colisiones e Implementación}

\begin{dsaquees}

Un estacionamiento 🅿️ donde una fórmula te asigna el cajón según tu
placa. El problema: a veces dos coches distintos reciben el MISMO cajón.
Eso es una colisión, y toda tabla hash necesita un plan para resolverla:

\begin{itemize}
\item Chaining: en cada cajón cabe una fila de coches (una lista). Si
      chocan, se forman detrás.
\item Open Addressing: si tu cajón está ocupado, buscas el siguiente
      libre siguiendo una regla.
\end{itemize}

Y si el estacionamiento se llena demasiado (Load Factor alto), se
construye uno más grande y se reacomodan todos los coches (rehashing).

\end{dsaquees}

\begin{dsaotraforma}

Una hash table promedia $O(1)$ SOLO si las colisiones se manejan bien.
Dos grandes estrategias:

\textbf{Chaining (encadenamiento)}: cada bucket guarda una lista
enlazada de pares que colisionaron. Simple y tolerante a Load Factors
altos. Es lo que usa \verb|std::unordered_map|.

\textbf{Open Addressing (direccionamiento abierto)}: todo vive en el
array; ante colisión se busca otro hueco mediante probing:
\begin{itemize}
\item Linear probing: prueba $i+1, i+2, \ldots$ (sufre "clustering").
\item Quadratic probing: prueba $i+1, i+4, i+9, \ldots$
\item Double hashing: el salto lo decide una segunda función hash.
\end{itemize}

El \textbf{Load Factor} $\alpha = n/m$ (elementos / buckets) mide qué
tan llena está la tabla. Cuando supera un umbral (~0.75), se hace
\textbf{rehashing}: se duplica el array y se reinsertan todas las claves.

\end{dsaotraforma}

\begin{dsadiagrama}
\begin{verbatim}
CHAINING (listas por bucket)      OPEN ADDRESSING (linear probing)

 [0] -> NULL                       hash("A")=2, hash("B")=2, hash("C")=3
 [1] -> NULL
 [2] -> (A) -> (B) -> (C)          [0] -
 [3] -> NULL                       [1] -
 [4] -> (D)                        [2] A     <- va a su hueco
                                   [3] B     <- ocupado el 2, prueba 3? no,
 Colisiones -> se apilan en la         C     <- desplazados en cadena
 lista del bucket.                 [4] -

 Load Factor a = n/m ; si a > 0.75 -> rehash (duplicar m y reinsertar).
\end{verbatim}
\end{dsadiagrama}

\begin{lstlisting}[language=C++]
IMPLEMENTACIÓN PROPIA (chaining)
1. Array de buckets; cada bucket es una lista de pares {clave, valor}.
2. índice = hash(clave) % capacidad.
3. Insert: si la clave existe en la lista, actualiza; si no, agrégala.
4. Get: recorre la lista del bucket buscando la clave.
5. Si load factor > umbral -> rehash (nueva capacidad, reinsertar todo).

📝 PSEUDOCÓDIGO (get con chaining)
──────────────────────────────────────────────────────────────

función get(clave):
    idx = hash(clave) % capacidad
    PARA cada par en buckets[idx]:
        SI par.clave == clave: return par.valor
    return NO_ENCONTRADO

💻 CÓDIGO C++ (hash map propio con chaining)
──────────────────────────────────────────────────────────────

#include <vector>
#include <list>
#include <utility>
using namespace std;

class MyHashMap {
    vector<list<pair<int,int>>> buckets;
    int capacity, size_;
    int idx(int key) { return (key % capacity + capacity) % capacity; }

    void rehash() {
        auto old = buckets;
        capacity *= 2;
        buckets.assign(capacity, {});
        size_ = 0;
        for (auto& b : old) for (auto& kv : b) put(kv.first, kv.second);
    }
public:
    MyHashMap(int cap = 8) : buckets(cap), capacity(cap), size_(0) {}

    void put(int key, int val) {
        auto& b = buckets[idx(key)];
        for (auto& kv : b) if (kv.first == key) { kv.second = val; return; }
        b.push_back({key, val});
        if (++size_ > capacity * 0.75) rehash();   // load factor
    }
    int get(int key) {
        for (auto& kv : buckets[idx(key)])
            if (kv.first == key) return kv.second;
        return -1;                                 // no encontrado
    }
    void remove(int key) {
        auto& b = buckets[idx(key)];
        for (auto it = b.begin(); it != b.end(); ++it)
            if (it->first == key) { b.erase(it); size_--; return; }
    }
};

⏱️ COMPLEJIDAD (Big-O)
──────────────────────────────────────────────────────────────

  Método            │ Promedio │ Peor caso
  ──────────────────┼──────────┼──────────
  Insert/Get/Delete │ O(1)     │ O(n)
  Rehash            │ O(n)     │ O(n) (amortizado O(1) por op)

* Peor caso O(n): función hash mala -> todo cae en un bucket.
\end{lstlisting}

\begin{dsaentrevistas}

✅ Chaining vs Open Addressing: chaining tolera load factors altos y
   borra fácil; open addressing usa mejor la cache y menos memoria por
   puntero, pero se degrada al llenarse. Sabe explicar ambos.

💡 Una buena hash function distribuye uniforme y es rápida. Para strings,
   el hash polinómico (\verb|h = h*31 + c|) es el clásico.

⚠️ Borrar en open addressing es delicado: no puedes dejar el hueco
   vacío (rompería las búsquedas por probing); se marca como "tombstone"
   (borrado lógico).

🔑 unordered\_map se degrada a O(n): con claves maliciosas que colisionan
   a propósito (hash flooding). Por eso algunos lenguajes aleatorizan la
   semilla del hash.

\end{dsaentrevistas}

\begin{dsaresumen}

• Colisión: dos claves caen en el mismo bucket. Toda tabla necesita plan.
• Chaining: lista por bucket (lo usa unordered\_map).
• Open Addressing: buscar otro hueco (linear/quadratic/double hashing).
• Load Factor $\alpha = n/m$; si pasa ~0.75 -> rehash (duplicar y reinsertar).
• Promedio O(1); peor caso O(n) con hash mala.
• Borrar en open addressing usa tombstones.
════════════════════════════════════════════════════════════════

\end{dsaresumen}
